\begin{lemma}
Let $1<p<\infty$, let $H$ be the extremal function of \noderef{ExtremalH}, let
$d=d(p)$ be the exponent of \noderef{DExp}, and let
$f\colon\mathbb{C}\to\mathbb{C}$ be a two-sided inverse of $H$, that is
\[
  f\bigl(H(\xi)\bigr)=\xi\quad\text{and}\quad H\bigl(f(z)\bigr)=z
  \qquad(\xi,z\in\mathbb{C}),
\]
which is the relation (55) of the paper.  Assume also that $f$ is of class
$C^\infty$ on $\mathbb{C}-\{0\}$.  Then
\begin{enumerate}
\item[(57)] $f(tz)=t^{d}f(z)$ for every $z\in\mathbb{C}$ and every real
$t\ge0$;
\item[(H)] for every word $w$ of length $l$, every real $t>0$ and every
$z\neq0$,
\[
  \bigl(D^{w}f\bigr)(tz)=t^{\,d-l}\,\bigl(D^{w}f\bigr)(z),
\]
where $D^{w}$ is the iterated partial derivative of \noderef{PartialWord}:
each partial derivative of $f$ of order $l$ is homogeneous of degree $d-l$.
\end{enumerate}

Clause (57) is the paper's (57) at paper.tex:597, and clause (H) is the
sentence that follows it, ``Hence, partial derivatives $D^\nu f$ are
homogeneous functions of degree $d-|\nu|$''.  The paper's displays (58) and
(59) are the two cases $l=k-1$ and $l=k+1$ of (H): since $d=k-1+\alpha$, the
exponents are $d-(k-1)=\alpha$ and $d-(k+1)=\alpha-2$, and (59) is the latter
rearranged.
\end{lemma}
\begin{proof}
Write $d=d(p)$.

\emph{Step 0: $d>0$.}  By \noderef{DExpEqGammaOneInv}, $\gamma_1(p)>0$ and
$\gamma_1(p)\,d=1$; hence $d=1/\gamma_1(p)>0$.

\emph{Step 1: proof of (57) for $t>0$.}  Let $t>0$ and $z\in\mathbb{C}$, and
put $\xi=f(z)$, so that $H(\xi)=z$ by the second inverse relation.  By the
second clause of \noderef{ExtremalPolarForm}, applied with this $t>0$ and this
$\xi$,
\[
  t\,H(\xi)=H\bigl(t^{d}\xi\bigr).
\]
The left-hand side is $tz$.  Applying $f$ to both sides and using the first
inverse relation at the point $t^{d}\xi$,
\[
  f(tz)=f\bigl(H(t^{d}\xi)\bigr)=t^{d}\xi=t^{d}f(z).
\]

\emph{Step 2: proof of (57) at $t=0$.}  By \noderef{ExtremalH}, $H(0)=0$; so
the first inverse relation at $\xi=0$ gives $f(0)=f(H(0))=0$.  Hence
$f(0\cdot z)=f(0)=0$, while $0^{d}f(z)=0$ because $0^{d}=0$ for $d>0$ by
Step 0.  So (57) holds at $t=0$ as well, and therefore for every $t\ge0$.

\emph{Step 3: how $D^{w}$ behaves under a dilation.}  Fix $t>0$.  We claim that
for every word $w$ and every function $F$ that is $C^\infty$ on
$\mathbb{C}-\{0\}$,
\begin{equation}\tag{1}
  D^{w}\bigl(\zeta\mapsto F(t\zeta)\bigr)(z)
  =t^{\,|w|}\bigl(D^{w}F\bigr)(tz)\qquad(z\neq0).
\end{equation}
We induct on $w$.  If $w$ is empty, $D^{w}$ is the identity by
\noderef{PartialWord} and both sides are $F(tz)=t^{0}F(tz)$.  For the step,
write $w=j\!:\!\!:w'$, so that $D^{w}=\partial_j\circ D^{w'}$ by
\noderef{PartialWord}, where $\partial_0$ and $\partial_1$ are the derivatives
in the directions $1$ and $i$.  The dilation $\zeta\mapsto t\zeta$ is an
$\mathbb{R}$-linear homeomorphism of $\mathbb{C}$ carrying $\mathbb{C}-\{0\}$
onto itself, with Fréchet derivative at every point equal to itself.  By the
inductive hypothesis the two functions
\[
  D^{w'}\bigl(\zeta\mapsto F(t\zeta)\bigr)
  \qquad\text{and}\qquad
  t^{\,|w'|}\bigl(D^{w'}F\bigr)(t\,\cdot)
\]
agree at every point of the open set $\mathbb{C}-\{0\}$, and both are
$C^\infty$ there, $F$ being $C^\infty$ on $\mathbb{C}-\{0\}$ and the dilation
mapping that set into itself.  Differentiating at $z\neq0$ in the direction
$v$ ($v=1$ if $j=0$, $v=i$ if $j=1$) and using the chain rule together with
$\mathbb{R}$-linearity of the derivative in its argument,
\[
  \partial_j\Bigl(t^{\,|w'|}\bigl(D^{w'}F\bigr)(t\,\cdot)\Bigr)(z)
  =t^{\,|w'|}\,\mathrm{d}\bigl(D^{w'}F\bigr)(tz)[t v]
  =t^{\,|w'|+1}\,\partial_j\bigl(D^{w'}F\bigr)(tz),
\]
which is (1) for $w$, since $|w|=|w'|+1$.

\emph{Step 4: proof of (H).}  Fix $t>0$ and a word $w$ of length $l$.  By
Step 1 the two functions $\zeta\mapsto f(t\zeta)$ and
$\zeta\mapsto t^{d}f(\zeta)$ agree at every point of $\mathbb{C}$, in
particular on the open set $\mathbb{C}-\{0\}$, where both are $C^\infty$ by
hypothesis and by Step 3's observation about the dilation.  Apply $D^{w}$ to
both and evaluate at $z\neq0$.  On the left, (1) with $F=f$ gives
$t^{\,l}\bigl(D^{w}f\bigr)(tz)$.  On the right, $t^{d}$ is a constant scalar
and $D^{w}$ is built from Fréchet derivatives, which are linear in the
function, so $D^{w}\bigl(t^{d}f\bigr)=t^{d}D^{w}f$ and the value is
$t^{d}\bigl(D^{w}f\bigr)(z)$.  Therefore
\[
  t^{\,l}\bigl(D^{w}f\bigr)(tz)=t^{d}\bigl(D^{w}f\bigr)(z),
\]
and dividing by $t^{\,l}>0$, together with $t^{d}/t^{\,l}=t^{\,d-l}$ for
$t>0$,
\[
  \bigl(D^{w}f\bigr)(tz)=t^{\,d-l}\bigl(D^{w}f\bigr)(z),
\]
which is (H).

\emph{Step 5: the two displays of the paper.}  Let $k$ and $\alpha$ be as in
(9), so that $d=k-1+\alpha$.  Taking $l=k-1$ in (H) gives
$\bigl(D^{w}f\bigr)(tz)=t^{\alpha}\bigl(D^{w}f\bigr)(z)$, which is (58); and
taking $l=k+1$ gives
$\bigl(D^{w}f\bigr)(tz)=t^{\alpha-2}\bigl(D^{w}f\bigr)(z)$, that is
$\bigl(D^{w}f\bigr)(z)=t^{\,2-\alpha}\bigl(D^{w}f\bigr)(tz)$, which is (59).
\end{proof}
